% tex.web §455's two font-dependent units. `em' is \fontdimen6 of the font in
% force and `ex' is \fontdimen5, so neither is a ratio to a point: the factor
% multiplies a number the FONT states, exactly as §453's factor multiplies an
% internal dimension. Nothing here selects a font, which is the point -- the
% current font after \input plain is cmr10, whose quad is 10.00002pt and whose
% x-height is 4.30554pt, and those are the numbers a document gets without
% asking for anything. A test that set its own font would prove the arithmetic
% and not the font.
% requires: cmr10 -- the numbers below are an INSTALLATION's cmr10.tfm, so a
% machine without one cannot be asked what they are. See tests/parity.rs.
\catcode`\{=1 \catcode`\}=2
\dimen0=1em \dimen1=1ex \dimen2=.5em \dimen3=2.5ex
\message{[\the\dimen0][\the\dimen1][\the\dimen2][\the\dimen3]}
% The sign rides on the whole product, and §107 truncates it toward zero.
\dimen0=-1.5ex \dimen1=0em \dimen2=3em
\message{[\the\dimen0][\the\dimen1][\number\dimen2]}
% §407's keyword scan: letters match either case, and a space in front of the
% unit is skipped.
\dimen0=1 em \dimen1=1EM \dimen2=1Ex
\message{[\the\dimen0][\the\dimen1][\the\dimen2]}
% All three components of a glue, and an infinite one beside a finite one --
% `fil' is five letters and `em' is two, so the longest match has to pick the
% right one of the two.
\skip0=1em plus .5em minus 1ex
\skip1=2ex plus 1fill
\message{[\the\skip0][\the\skip1]}
% A quad is a dimension like any other once it is scanned: it compares, it
% advances, and it coerces to its count of scaled points where a number is
% wanted (§430).
\dimen0=1em \advance\dimen0 by 1ex \multiply\dimen0 by 2
\count1=\dimen0
\message{[\the\dimen0][\the\count1]}
\ifdim 1em>1ex \message{[quad wider]}\else\message{[no]}\fi
% §455 runs before §453 reaches `true', so the letters after it are matched
% against the PHYSICAL units only and `em' is not one of them. Scanning the
% font units after the keyword would accept this line.
\dimen0=1true pt
\message{[\the\dimen0]}
\end
